% !TeX root = ../../Book5.tex
% 纲要标题：### 6.4 不可约分类
% 纲要位置：工作记录/计划纲要.md，第 4245 行。
% * 特征零下的不可约性
% * 不同分拆给出不同不可约表示
% * 分类完备性
% * 调用第6.3节已建立的 Garnir 关系及标准基，在本节依次完成不可约性、非同构性与完备性；不使用尚未证明的钩长公式或 Frobenius 特征映射
% * 用 Specht 模上的不变配对与子模定理证明特征零不可约性，明确正特征下哪一步失效
% * 用支配序和 Young 子群分析不同分拆的非同构性；不把参数个数相等当作非同构性的证明
% * 在不可约性和非同构性建立后，结合共轭类个数与半单群代数完成分类完备性；处理空分拆，并用整数标准基把有理数上的作用同态扩张到任意特征零域
% * 分类之后分别讨论 Young 对称化子的归一化与幂等元、正特征中的配对退化与可约性、Young 置换模分解的支配三角性及维数平方和的核验；正式证明特征二下形状 \((2,2)\) 的 Specht 模配对为零但仍绝对不可约，使标题区分新构造、反例、推论与事后检验
\section{不可约分类}
\label{b5:sec:0604}

标准基解决的是向量空间的大小，尚未说明群作用能否保持一个更小的非零子空间。本节在表类基上引入不变配对；列反对称化对这个配对有一种特别简单的秩一作用，从而控制 Specht 模的所有子模。

\subsection{特征零下的不可约性}
\label{b5:subsec:0604:01}

在 \(M_R^\lambda\) 上把不同行表类规定为正交，并令每个基向量与自身的配对为 \(1\)，得到对称双线性型
\[
	B\left(\sum_u a_u\{u\},\sum_u d_u\{u\}\right)=\sum_u a_ud_u.
\]
群作用置换这组基，所以 \(B(gv,gw)=B(v,w)\)。这里采用双线性型；在复数域上不把它与共轭线性的 Hermitian 内积混为一谈。

\begin{lemma}\label{b5:lem:specht-rank-one}
	设 \(R\) 为交换含幺环，\(n\) 为非负整数，\(\lambda\vdash n\)，\(t\) 为形状 \(\lambda\) 的 Young 表，\(C_t\) 为其列群。记 \(M_R^\lambda\) 为行表类置换模，\(b_t=\sum_{c\in C_t}\operatorname{sgn}(c)c\)、\(e_t=b_t\{t\}\)，\(B\) 为以行表类基为标准正交基的对称双线性型。对任意 \(v\in M_R^\lambda\)，有
	\begin{equation}\label{b5:eq:specht-rank-one}
		b_tv=B(v,e_t)e_t.
	\end{equation}
\end{lemma}
\begin{proof}
	先令 \(v=\{u\}\)。若 \(t\) 的某一列有两个数字位于 \(u\) 的同一行，交换它们的换位既属于 \(C_t\) 又固定 \(\{u\}\)；把列群求和配成符号相反的二项，即得 \(b_t\{u\}=0\)。

	若不存在这样的两个数字，则 \(u\) 的第一行共有 \(\lambda_1\) 个数字，而 \(t\) 恰有 \(\lambda_1\) 列，每列至多贡献一个；所以每列都恰贡献一个。可用各列内部的置换把这些数字送到 \(t\) 的第一行。删去已匹配的第一行后，尚非空的列数恰为 \(\lambda_2\)，而第二行也有 \(\lambda_2\) 个数字，每列仍至多贡献一个，故再次每列恰贡献一个。以后各行同理。逐行进行，得到 \(c\in C_t\) 使 \(\{u\}=\{ct\}\)。于是 \(b_t\{u\}=\operatorname{sgn}(c)e_t\)。这证明 \(b_tv\) 总是 \(e_t\) 的标量倍数。

	因为 \(c\) 与 \(c^{-1}\) 符号相同，\(b_t\) 关于 \(B\) 自伴随。故 \(b_tv\) 中 \(\{t\}\) 的系数为 \(B\bigl(b_tv,\{t\}\bigr)=B\bigl(v,b_t\{t\}\bigr)=B(v,e_t)\)。而 \(e_t\) 中 \(\{t\}\) 的系数为 \(1\)，所求标量即为此数。
\end{proof}

例如对任意交换含幺环 \(R\)，在\cref{b5:exm:specht-21}的形状 \((2,1)\) 中取 \(t=12|3\)，以 \(v_i\) 表示第二行是 \(i\) 的表类。此时
\[
	b_t=1-(13),\qquad e_t=v_3-v_1,
\]
而置换 \((13)\) 交换 \(v_1,v_3\) 并固定 \(v_2\)，所以
\[
	b_tv_1=-e_t,\qquad b_tv_2=0,\qquad b_tv_3=e_t.
\]
对 \(v=a_1v_1+a_2v_2+a_3v_3\)，便有
\[
	b_tv=(a_3-a_1)e_t=B(v,e_t)e_t.
\]
列反对称化把表类置换模送到 \(Re_t\)；配对 \(B(v,e_t)\) 测出其像在 \(e_t\) 前的系数。当底环是域时，只要子模中某个向量的这个系数非零，就能从它取出 \(e_t\)，再由群作用得到所有列交错元。这给出以下子模定理。

\begin{theorem}[子模定理]\label{b5:thm:specht-submodule}
	设 \(k\) 为任意域，\(n\) 为非负整数，\(\lambda\vdash n\)，\(M_k^\lambda\) 为行表类置换模，\(S_k^\lambda\subseteq M_k^\lambda\) 为 Specht 模，\(U\) 为 \(M_k^\lambda\) 的左 \(k[S_n]\)-子模。则至少有下列一项成立：
	\[
		S_k^\lambda\subseteq U,\qquad U\subseteq(S_k^\lambda)^\perp,
	\]
	其中正交补取在 \(M_k^\lambda\) 中，\(B\) 为以行表类基为标准正交基的对称双线性型。
\end{theorem}
\begin{proof}
	若第二项不成立，因为全部 \(e_t\) 生成 \(S_k^\lambda\)，存在 \(u\in U\) 与某个 \(t\) 使 \(B(u,e_t)\ne0\)。由\eqref{b5:eq:specht-rank-one}， \(b_tu=B(u,e_t)e_t\in U\)。域中非零标量可逆，故 \(e_t\in U\)；这个向量的 \(S_n\)-轨道生成整个 Specht 模，所以第一项成立。
\end{proof}

\begin{proposition}\label{b5:prop:specht-char-zero-simple}
	设 \(k\) 为特征零域，\(n\) 为非负整数，\(\lambda\vdash n\)，\(M_k^\lambda\) 为行表类置换模，\(S_k^\lambda\) 为其中的 Specht 模。以 \(M_k^\lambda\) 的行表类基为标准正交基的对称双线性型 \(B\)，在 \(S_k^\lambda\) 上非退化，且 \(S_k^\lambda\) 作为左 \(k[S_n]\)-模绝对不可约。
\end{proposition}
\begin{proof}
	在整数标准基中写出 \(B\) 的 Gram 矩阵 \(G_\lambda\)，其条目为整数。把它看作实矩阵时，由\cref{b5:thm:specht-standard-basis}，这些基向量在实表类空间中线性无关，而表类配对为正定型。因此其 Gram 矩阵为正定实对称矩阵，特征值全为正，故 \(\det G_\lambda>0\)。又因条目为整数，这个行列式是非零整数。它在任意特征零域中的像仍非零，故限制型非退化。

	若 \(0\ne U\subseteq S_k^\lambda\) 是子模，子模定理或者给出 \(S_k^\lambda\subseteq U\)，或者给出 \(U\subseteq S_k^\lambda\cap(S_k^\lambda)^\perp=0\)。后者不可能，故 \(U=S_k^\lambda\)。对 \(k\) 的任何扩域，标准基与同一个非零 Gram 行列式仍成立，所以上述不可约性在扩域后保持，即为绝对不可约。
\end{proof}

\subsection{不同分拆与非同构表示}
\label{b5:subsec:0604:02}

\begin{proposition}\label{b5:prop:specht-dominance}
	设 \(k\) 为特征零域，\(n\) 为非负整数，\(\lambda,\mu\vdash n\)，\(S_k^\lambda,S_k^\mu\) 为相应 Specht 左 \(k[S_n]\)-模，\(M_k^\lambda,M_k^\mu\) 为相应行表类置换模。若 \(\operatorname{Hom}_{k[S_n]}(S_k^\lambda,M_k^\mu)\ne0\)，则 \(\lambda\trianglerighteq\mu\)。此外，
	\[
		\dim_k\operatorname{Hom}_{k[S_n]}(S_k^\lambda,M_k^\lambda)=1.
	\]
	特别地，\(S_k^\lambda\cong S_k^\mu\) 当且仅当 \(\lambda=\mu\)。
\end{proposition}
\begin{proof}
	固定一个 \(\lambda\)-表 \(t\)，并令 \(\varphi:S_k^\lambda\rightarrow M_k^\mu\) 非零。因为 \(e_t\) 生成整个模，\(\varphi(e_t)\ne0\)。列群对 \(e_t\) 按符号作用，因此
	\[
		b_t\cdot\varphi(e_t)=\varphi(b_te_t)=|C_t|\cdot\varphi(e_t)\ne0.
	\]
	所以至少有一个 \(\mu\)-行表类 \(\{u\}\) 满足 \(b_t\{u\}\ne0\)。与秩一引理中的配对消去相同，\(u\) 的任何一行不能含有 \(t\) 同一列的两个数字。前 \(a\) 行从第 \(j\) 列至多取 \(\min(a,\lambda'_j)\) 个数字，故
	\[
		\mu_1+\cdots+\mu_a \leq\sum_j\min(a,\lambda'_j) =\lambda_1+\cdots+\lambda_a.
	\]
	这就是支配条件。

	当 \(\mu=\lambda\) 时，\cref{b5:lem:specht-rank-one}又说明 \(b_t\cdot\varphi(e_t)\in ke_t\)。因 \(|C_t|\) 可逆，\(\varphi(e_t)\) 为 \(e_t\) 的标量倍数；由生成性，\(\varphi\) 是包含映射的标量倍数。包含映射确实非零，故 Hom 空间恰为一维。

	最后，若两个 Specht 模同构，把同构与 \(S_k^\mu\subseteq M_k^\mu\) 复合，得到 \(\lambda\trianglerighteq\mu\)；用逆同构得到 \(\mu\trianglerighteq\lambda\)。支配序的反对称性给出 \(\lambda=\mu\)。反方向显然成立。
\end{proof}

证明中的计数描述了一个行列之间的容量限制。一个 \(\mu\)-表类要使 \(b_t\{u\}\ne0\)，每一行至多从 \(t\) 的每一列取一个数字；否则交换同一列中落到同一行的两个数字，列交错求和就成对相消。对前 \(a\) 行合起来，第 \(j\) 列既只能提供 \(\lambda'_j\) 个数字，又至多提供 \(a\) 个，所以其容量为 \(\min(a,\lambda'_j)\)。总容量恰是 \(\lambda\) 的前 \(a\) 行方格数，因此必须容得下 \(\mu\) 的前 \(a\) 行。

\ProseParagraphBreak
例如 \(\lambda=(2,2)\) 只有两列，而 \(\mu=(3,1)\) 的第一行需要三个数字。无论怎样安排，这三个数字中都有两个来自 \(t\) 的同一列，因此每个 \((3,1)\)-表类都被 \(b_t\) 零化。第一行的不等式已经失败：\(3>\min(1,2)+\min(1,2)=2\)。在特征零下，若存在非零同态 \(S_k^{(2,2)}\rightarrow M_k^{(3,1)}\)，它在生成元上的像又必须满足 \(b_t\varphi(e_t)=4\varphi(e_t)\ne0\)，与这个相消计算矛盾。

\subsection{分类完备性}
\label{b5:subsec:0604:03}

\begin{theorem}[对称群的不可约分类]\label{b5:thm:specht-classification}
	设 \(k\) 为特征零域，\(n\geq0\) 为整数。全部 Specht 模 \(S_k^\lambda\)，其中 \(\lambda\vdash n\)，是互不同构的绝对不可约左 \(k[S_n]\)-模，并穷尽所有不可约有限维左 \(k[S_n]\)-模。
\end{theorem}
\begin{proof}
	当 \(n=0\) 时，\(S_0\) 是平凡群，唯一的空分拆对应 Specht 模 \(k\)，结论成立。以下设 \(n\geq1\)。不可约性和非同构性已经分别由\cref{b5:prop:specht-char-zero-simple}及\cref{b5:prop:specht-dominance}证明。先取 \(k=\mathbb C\)。\cref{b5:prop:symmetric-conjugacy-partitions}说明共轭类数等于分拆数，而\cref{b5:thm:irreducible-character-basis}说明不可约复表示的类型数等于共轭类数。我们已经为每个分拆构造出互不同构的不可约表示，因此没有其他复不可约类型。

	为把完备性推广到任意特征零域，先利用整数模型。\(S_{\mathbb Z}^\lambda\) 被群作用保持，标准基定理又说明它有整数标准基，所以每个群元素的作用矩阵都有整数条目。用这些矩阵得到一个在有理数上定义的代数同态
	\[
		\Phi_{\mathbb Q}:\mathbb Q[S_n]\longrightarrow \prod_{\lambda\vdash n} \operatorname{End}_{\mathbb Q}(S_{\mathbb Q}^\lambda).
	\]
	对任意特征零域 \(k\)，标准基定理给出
	\[
		k\otimes_{\mathbb Q}S_{\mathbb Q}^\lambda \xrightarrow{\ \sim\ }S_k^\lambda, \qquad a\otimes e_t\longmapsto ae_t.
	\]
	两边的群作用矩阵都由同一组整数矩阵得到，故此同构等变。在这些基下，\(\operatorname{End}_{\mathbb Q}(S_{\mathbb Q}^\lambda)\) 的矩阵单位基也扩张为 \(\operatorname{End}_k(S_k^\lambda)\) 的基。因此 \(k\) 上的作用同态 \(\Phi_k\) 正是 \(k\otimes_{\mathbb Q}\Phi_{\mathbb Q}\)，其矩阵由 \(\Phi_{\mathbb Q}\) 的矩阵逐项扩张系数得到。

	扩张到复数后，此映射单射：若群代数元素在所有不可约上作用为零，由\cref{b5:thm:maschke}，复正则表示是这些不可约表示的直和，故该元素在正则表示上也作用为零；作用于单位元即知该元素为零。由\cref{b5:thm:row-column-orthogonality}在单位共轭类处得到的维数平方和，目标维数为 \(\sum_\lambda(f^\lambda)^2=n!\)，所以这个单射还是同构。

	因此 \(\Phi_{\mathbb Q}\) 在所选基中是方阵，其有理行列式非零。在任意特征零域 \(k\) 中，这个行列式的像仍非零，故
	\begin{equation}\label{b5:eq:symmetric-split-algebra}
		k[S_n]\cong\prod_{\lambda\vdash n}M_{f^\lambda}(k).
	\end{equation}
	在这个直积中，令 \(z_\lambda\) 在形状 \(\lambda\) 的矩阵块上为单位矩阵、在其他块上为零。它们是一组正交中心幂等元，且 \(1=\sum_\lambda z_\lambda\)。任一模 \(M\) 因而分解为 \(\bigoplus_\lambda z_\lambda M\)；若 \(M\) 不可约，其中恰有一项非零，所以只属于一个矩阵块。设 \(V\) 是矩阵代数 \(M_d(k)\) 的不可约模。取 \(0\ne v\in V\)，由 \(1=\sum_iE_{ii}\) 可选 \(i\) 使 \(E_{ii}v\ne0\)。记 \(e_1,\ldots,e_d\) 为 \(k^d\) 的标准列向量基，定义
	\[
		\psi:k^d\longrightarrow V, \qquad e_j\longmapsto E_{ji}v.
	\]
	矩阵单位关系 \(E_{ab}E_{ji}=\delta_{bj}E_{ai}\) 说明 \(\psi\) 是模同态，而 \(\psi(e_i)=E_{ii}v\ne0\)。列向量模 \(k^d\) 简单，所以这个非零同态的核为零；其像是不可约模 \(V\) 的非零子模，故等于 \(V\)。于是 \(\psi\) 为同构。这证明\eqref{b5:eq:symmetric-split-algebra}所列的 Specht 模也穷尽 \(k\) 上全部类型。
\end{proof}

每个分拆已给出一个表示。证明它们各自不可约且彼此不同后，分拆数等于共轭类数就保证全部不可约类型均已出现，从而完成分类。

\ProseParagraphBreak
从复数推广到任意特征零域，依靠的是对称群表示的一个具体性质：每个不可约复表示都是某个 \(S_{\mathbb Z}^\lambda\) 的复化，群作用在标准基中由整数矩阵给出。这使全部不可约表示同时有同维的有理数域模型，才能先构造 \(\Phi_{\mathbb Q}\)，再由复扩张检验它的有理行列式非零。一般有限群并不自动具有这样的模型。\Cref{b5:prop:rational-values-no-real-model}中的 \(Q_8\) 二维不可约复特征标虽然取整数值，却不能在 \(\mathbb Q\) 上实现为二维表示；仅知道复不可约类型及其特征标值，并不能得出有理群代数全部分裂。

\subsection{Young 对称化子的归一化与幂等元}
\label{b5:subsec:0604:04}

标准基使 Gram 行列式、基域扩张与群作用矩阵都有整系数的共同来源。它也使 Young 对称化子可以精确归一化：

\begin{proposition}\label{b5:prop:young-quasi-idempotent}
	设 \(k\) 为特征零域，\(n\) 为非负整数，\(\lambda\vdash n\)，\(t\) 为形状 \(\lambda\) 的 Young 表，\(R_t,C_t\) 分别为其行群、列群。令 \(a_t=\sum_{r\in R_t}r\)、\(b_t=\sum_{c\in C_t}\operatorname{sgn}(c)c\)、\(c_t=b_ta_t\in k[S_n]\)，并记 \(f^\lambda\) 为形状 \(\lambda\) 的标准表数。则
	\[
		c_t^2=\frac{n!}{f^\lambda}c_t,
	\]
	且 \(\dfrac{f^\lambda}{n!}c_t\) 是一个幂等元，它生成的左理想同构于 Specht 模 \(S_k^\lambda\)。
\end{proposition}
\begin{proof}
	由\cref{b5:prop:young-ideal-model}，计算 \(c_t^2\) 等价于计算 \(b_ta_te_t\)。秩一公式给出 \(b_ta_te_t=B(a_te_t,e_t)e_t\)。记这个整数系数为 \(\alpha\)，便有 \(c_t^2=\alpha c_t\)。

	在实表类空间中，\(a_t\) 自伴随且 \(a_t^2=|R_t|a_t\)，所以
	\[
		\alpha=\frac{B(a_te_t,a_te_t)}{|R_t|}>0.
	\]
	严格不等号的理由是 \(a_te_t\) 中 \(\{t\}\) 的系数为 \(|R_t|\)：每个行置换都固定 \(\{t\}\)，而 \(e_t\) 在该处系数为 \(1\)。因此 \(\alpha\ne0\)。

	右乘 \(c_t/\alpha\) 是 \(k[S_n]\) 上的投影，像是 \(k[S_n]c_t\)，其维数为 \(f^\lambda\)。由 \(R_t\cap C_t=\{1\}\)，\(c_t\) 中单位元系数为 \(1\)，所以在群元素基下，右乘 \(c_t/\alpha\) 的每个对角元都是 \(1/\alpha\)，迹为 \(n!/\alpha\)。投影的迹等于像的维数，得到 \(\alpha=n!/f^\lambda\)。
\end{proof}

这项构造可与 Etingof 等人的 Young 对称化子证明对照\cite[Sections 5.12--5.13, pp. 110--113]{EtingofEtAl2011RepresentationTheory}。该书采用预先归一化的行、列求和及相反乘法次序，因而相应的倍数也随约定改变；本节的归一化是在当前模型内算出的。

\ProseParagraphBreak
令 \(d=f^\lambda\)，并记 \(e=(d/n!)c_t\)。这个幂等元与\eqref{b5:eq:central-idempotent}中的本原中心幂等元承担不同的作用。在\eqref{b5:eq:symmetric-split-algebra}的分解下，形状 \(\lambda\) 对应矩阵块 \(M_d(k)\)，令 \(z_\lambda\) 为在此块上等于 \(I_d\)、在其他块上等于零的本原中心幂等元。左理想 \(k[S_n]z_\lambda\) 是整个矩阵块，维数为 \(d^2\)；而 \(k[S_n]e\cong S_k^\lambda\) 是一个简单左理想，维数只有 \(d\)。

\ProseParagraphBreak
具体地，\(k[S_n]e\) 的简单性说明 \(e\) 只在形状 \(\lambda\) 的块中非零。该块中的幂等矩阵若秩为 \(r\)，所生成的左理想维数为 \(dr\)，所以这里 \(r=1\)；选取像与核的基后，\(e\) 在该块中成为 \(E_{11}\)。当 \(d>1\) 时，它不是中心元素，而 \(z_\lambda\) 对应 \(I_d\)。例如 \(S_3\) 的形状 \((2,1)\) 有 \(d=2\)，故 \(e=c_t/3\) 生成二维简单左理想，\(z_{(2,1)}\) 则生成四维矩阵块。按列分解 \(M_d(k)=\bigoplus_{i=1}^dM_d(k)E_{ii}\)，也就看到相应不可约模在正则表示中出现 \(d\) 次。

\subsection{正特征中的配对退化与可约性}
\label{b5:subsec:0604:05}

子模定理本身允许任意特征；不可约性证明需要的额外事实是 Specht 模上配对非退化。非零整数 Gram 行列式在正特征中可能变成零，这正是不能照搬证明的位置。

\ProseParagraphBreak
非同构性的证明还有另一处特征条件。\Cref{b5:prop:specht-dominance}用 \(b_t\varphi(e_t)=|C_t|\varphi(e_t)\) 推出左边非零，这要求列群阶在域中非零。若特征 \(p\) 整除 \(|C_t|\)，即使 \(\varphi(e_t)\ne0\)，也不能作此推断。因此行列容量计数仍成立，却未必能从非零同态得到一个未被 \(b_t\) 零化的表类。最小的例子是特征二域 \(k\) 上的 \(S_2\)：由\cref{b5:prop:specht-column-sign}，\(S_k^{(2)}\) 与 \(S_k^{(1,1)}\) 分别是平凡表示与符号表示，但 \(-1=1\) 使它们同构。配对的退化与列群阶变成零，是前面两段论证中分别需要辨认的条件。

\begin{example}\label{b5:exm:specht-mod-three}
	设 \(k\) 为特征 \(3\) 的域。由\cref{b5:exm:specht-21}， \(S_k^{(2,1)}\) 是三点置换模中的坐标和为零子空间。向量 \(v_1+v_2+v_3\) 非零，因 \(3=0\) 而属于这个二维子空间，并生成一维平凡子模，所以 \(S_k^{(2,1)}\) 可约。在基 \(v_2-v_1,v_3-v_1\) 中，
	\[
		G_{(2,1)}=
		\begin{pmatrix}2&1\\1&2\end{pmatrix},
		\qquad \det G_{(2,1)}=3.
	\]
	其根恰好是上述平凡子模。标准基仍有两个向量，退化的是配对与不可约性。
\end{example}

配对退化并不必然产生非零真子模。形状 \((2,2)\) 在特征二时，限制配对甚至处处为零，群作用却仍绝对不可约。

\begin{proposition}\label{b5:prop:specht-22-char-two}
	设 \(k\) 为域，\(t_1=12|34\)、\(t_2=13|24\) 为形状 \((2,2)\) 的两个标准 Young 表，并令 \(e_i=e_{t_i}\)。这两个整数列交错元给出 \(S_{\mathbb Z}^{(2,2)}\) 的标准基，换系数后给出 \(S_k^{(2,2)}\) 的基。以行表类基为标准正交基的对称双线性型在此基中具有 Gram 矩阵
	\[
		G_{(2,2)}=\begin{pmatrix}4&2\\2&4\end{pmatrix}.
	\]
	在同一基中，作用矩阵为
	\[
		\rho((12))=\rho((34))=\begin{pmatrix}1&0\\-1&-1\end{pmatrix},
		\qquad
		\rho((23))=\begin{pmatrix}0&1\\1&0\end{pmatrix}.
	\]
	这些矩阵先在整数上成立，再将条目映入 \(k\)。 若 \(\operatorname{char}k=2\)，该限制型为零，但 \(S_k^{(2,2)}\) 仍绝对不可约。
\end{proposition}
\begin{proof}
	先在整数表类模中展开。两个表的列群分别由 \((13),(24)\) 及 \((12),(34)\) 生成，故
	\[
		\begin{aligned}
			e_1&=\{12|34\}-\{23|14\}-\{14|23\}+\{34|12\},\\
			e_2&=\{13|24\}-\{23|14\}-\{14|23\}+\{24|13\}.
		\end{aligned}
	\]
	每个向量有四项系数为 \(\pm1\)，两者共有的两项系数均为 \(-1\)，因此得到所列 Gram 矩阵。置换这些表类又给出
	\[
		\begin{aligned}
			(12)e_1&=e_1-e_2,&\qquad (12)e_2&=-e_2,\\
			(34)e_1&=e_1-e_2,& (34)e_2&=-e_2,\\
			(23)e_1&=e_2,& (23)e_2&=e_1.
		\end{aligned}
	\]
	采用作用矩阵的各列记录基向量像的约定，即得所列作用矩阵。这些整数等式换系数后仍成立。

	现在设 \(\operatorname{char}k=2\)，两个矩阵分别成为
	\[
		A=\begin{pmatrix}1&0\\1&1\end{pmatrix},\qquad
		B=\begin{pmatrix}0&1\\1&0\end{pmatrix}.
	\]
	二维空间的非零真子空间只能是一条直线。若直线被 \(A\) 保持，\(A\) 在其上按某个标量作用；由 \((A-I)^2=0\)，这个标量必须是 \(1\)，所以该直线包含在 \(\ker(A-I)=ke_2\) 中，只能等于 \(ke_2\)。但 \(Be_2=e_1\)，它不被 \(B\) 保持。因此不存在非零真子表示。对任意扩域 \(K/k\)，标准基仍是 \(e_1,e_2\)，两个矩阵也相同，上述论证仍排除 \(K\)-不变直线，故表示绝对不可约。另一方面，Gram 矩阵的所有条目在特征二中均为零。
\end{proof}

这里 \(S_k^{(2,2)}\subseteq(S_k^{(2,2)})^\perp\)，所以子模定理的正交补分支不能排除真子模，须直接考察群作用。在\subsecref{b5:subsec:0606:03}中，忘掉两行编号所得的无序配对商映射，其标准列交错元的像带有系数 \(2\)，因而在特征二时对整个 Specht 模为零。这说明某个交织映射在换系数后可以失效，而它的源模仍是非零的绝对不可约表示。

\subsection{Young 置换模分解的支配三角性}
\label{b5:subsec:0604:06}

设 \(k\) 为特征零域。由\eqref{b5:eq:symmetric-split-algebra}，\(S_k^\lambda\) 上的群代数作用为全部矩阵代数，所以与群作用对易的自同态只有标量：与各 \(E_{ii}\) 交换使其矩阵为对角矩阵，再与各 \(E_{ij}\) 交换，使所有对角元相等。\Cref{b5:thm:maschke}又使 \(M_k^\mu\) 完全可约，因而 \(S_k^\lambda\) 在其中的重数既等于 \(\dim_k\operatorname{Hom}_{k[S_n]}(S_k^\lambda,M_k^\mu)\)，也等于 \(\dim_k\operatorname{Hom}_{k[S_n]}(M_k^\mu,S_k^\lambda)\)。由于 \(M_k^\mu\cong\operatorname{Ind}_{S_\mu}^{S_n}k\)，将\cref{b5:thm:frobenius-reciprocity}用于平凡 \(k[S_\mu]\)-模 \(k\)，后一 Hom 空间同构于
\[
	\operatorname{Hom}_{k[S_\mu]}(k,\operatorname{Res}_{S_\mu}^{S_n}S_k^\lambda) \cong (S_k^\lambda)^{S_\mu},
\]
最后一个同构把同态送到 \(1\) 的像。这便把重数与 Young 子群不动向量联系起来，而\cref{b5:prop:specht-dominance}对前一 Hom 空间的限制给出三角性：
\[
	M_k^\mu\cong S_k^\mu\oplus \bigoplus_{\substack{\lambda\trianglerighteq\mu\\\lambda\ne\mu}} (S_k^\lambda)^{\oplus m_{\lambda\mu}}, \qquad m_{\lambda\mu}\in\mathbb Z_{\geq0}.
\]
这里尚未计算所有重数，但对角重数已经精确等于 \(1\)。后面建立 Frobenius 特征映射时，正是这项有方向的三角性帮助识别 Schur 函数。

\subsection{维数平方和的核验}
\label{b5:subsec:0604:07}

分类还给出一个事后检验维数的方法：
\[
	n!=\sum_{\lambda\vdash n}(f^\lambda)^2.
\]
左边来自群的阶，右边来自各个标准表集合的大小。这个恒等式在分类完成时已经成立；钩长公式将使每一项都能迅速算出。它适合检查是否漏了分拆或算错维数，却不能在没有不可约性和非同构性证明时单独充当分类论证。
